% 第 IX 讲，第 9 节；原 PDF 427—447 页。
\SourcePage{427}
\SGAsection{9}{单值性配对 $\underline M_\ell\otimes\underline M'_\ell\longrightarrow\mathbb Z_{\ell S}$ 的定义}
\textbf{9.1。}本节假定 $S$ 是\Term{henselian}{Hensel}的，并且 $A_K$ 在 $S$ 上具有\Term{semistable-reduction}{半稳定约化}；\Term{equivalence}{等价}地，$A'_K$ 在 $S$ 上具有\Term{semistable-reduction}{半稳定约化}（3.5.1）。于是有 $S$ 上的两个\Term{twisted-constant-module}{扭常值模} $\underline M$ 和 $\underline M'$（5.5.2）。在 $S$ 的某个适当的\Term{finite}{有限}\Term{etale}{平展}\Term{covering}{覆盖}上，它们成为取值分别为 $M$ 和 $M'$ 的\Term{constant-group}{常值群}；若 $S$ \Term{strictly-local}{严格局部}，即 $k$ \Term{separably-closed}{可分闭}，则在 $S$ 本身上已经如此。$M$ 和 $M'$ 都（非\Term{canonical}{典范}地）\Term{isomorphism}{同构}于 $\mathbb Z^r$，其中 $r$ 为 $A_0$ 与 $A'_0$ 的最大\Term{torus}{环面}的共同\Term{dimension}{维数}。对于每个\Term{prime-number}{素数} $\ell$，我们考察相应的 $\ell$ \Term{adic}{进}\Term{sheaf}{层}
\begin{equation}\tag{9.1.1}
 \underline M_\ell=\underline M\otimes_{\mathbb Z}\mathbb Z_\ell,
 \qquad \underline M'_\ell=\underline M'\otimes_{\mathbb Z}\mathbb Z_\ell,
\end{equation}
其几何意义由公式 (5.5.7) 与 (5.5.8) 阐明。本节将定义一个\Term{canonical}{典范}\Term{pairing}{配对}（称为“\Term{monodromy-pairing}{单值性配对}”）
\begin{equation}\tag{9.1.2}
 u_\ell:\underline M_\ell\otimes\underline M'_\ell\longrightarrow\mathbb Z_\ell.
\end{equation}
它将用 Barsotti--Tate 的\Term{pro-ell-group}{pro-$\ell$ 群} $T_\ell(A^0)$ 与 $T_\ell(A'^0)$ 来定义，甚至只用前者的\Term{generic-fiber}{泛纤维} $T_\ell(A_K)$ 就能定义（当 $\ell=p$ 时，须能应用 Tate \Term{theorem}{定理} [33, th.~4]，即目前须处于\Term{mixed-characteristic}{混合特征}情形）。当 $\ell\ne p$ 时，也就是说，\Term{pairing}{配对} (9.1.2) 可用 (2.5.3) 中考察过的 \Term{galois-module}{Galois $\pi$-模}
\begin{equation}\tag{9.1.3}
 U=T_\ell(A_K)(\overline K)=T_\ell(A_{\overline K})
\end{equation}
来描述。
\SourcePage{428}
更确切地说，我们将看到，知道这个\Term{pairing}{配对}本质上\Term{equivalence}{等价}于知道\Term{monodromy-group}{单值群} $I$ 在 $U$ 上的\Term{action}{作用}。

\textbf{9.2。$\ell\ne p$ 的情形。}首先考虑较简单的 $\ell\ne p$ 情形，并沿用 (2.5.3) 引入的记号。注意，忽略挠部分时，关系 $V=U^I$ 也意味着 $V^\perp$ 是由所有元素 $gx-x$（$x\in U$，$g\in I$）生成的\Term{module}{模}；无论如何，总有
\begin{equation}\tag{9.2.1}
 gx-x\in V^\perp\qquad(x\in U,\ g\in I).
\end{equation}
因此，对 $g\in I$，有
\begin{equation}\tag{9.2.2}
 gx=x+u(g).x,
\end{equation}
其中 $g\mapsto u(g)$ 是从 $U$ 到 $V^\perp$ 的一个\Term{homomorphism}{同态}，在 $V=U^I$ 上为零；\Term{equivalence}{等价}地，
\begin{equation}\tag{9.2.3}
 u\in\operatorname{Hom}(U/V,V^\perp)\simeq(U/V)^\vee\otimes V^\perp.
\end{equation}
在所考察的\Term{semistable-reduction}{半稳定约化}情形，即 $V^\perp\subset V$ 时，公式 (9.2.2) 表明 $g$ 在 $U$ 中的\Term{action}{作用}是一个平行于 $V^\perp=W$ 的\Term{transvection}{剪切变换}（transvection）；这些\Term{transvection}{剪切变换}的复合，正是相应\Term{homomorphism}{同态} $u$ 的相加 (9.2.3)。所以\Term{inertia-group}{惯性群} $I$ 在 $U$ 上的\Term{action}{作用}，完全由一个群\Term{homomorphism}{同态}
\begin{equation}\tag{9.2.4}
 I\longrightarrow(U/V)^\vee\otimes V^\perp
\end{equation}
来描述。
\SourcePage{429}
将\Term{orthogonality-theorem}{正交性定理}同时用于 $A_K$ 与 $A'_K$，得到
\[
 V^\perp=W,\qquad (U/V)^\vee\simeq W'(-1),
\]
其中 $W'$ 是 $U'=T_\ell(A'(\overline K))$ 的\Term{toric-part}{环面部分}，而 $(-1)$ 表示 \Term{tate-twist}{Tate 扭转}。另一方面，\Term{homomorphism}{同态} (9.2.4) 通过 $I(\ell)\simeq\mathbb Z_\ell(1)$ \Term{factorization}{分解}，故 (9.2.4) 可以解释为一个\Term{canonical}{典范}\Term{homomorphism}{同态}
\begin{equation}\tag{9.2.5}
 u:\mathbb Z_\ell(1)\longrightarrow W'\otimes W(-1),
\end{equation}
也可以解释为一个元素
\begin{equation}\tag{9.2.6}
 u\in W'\otimes W(-2).
\end{equation}
回顾 $W$ 与 $W'$ 的定义：它们是两个 $\ell$ \Term{adic}{进}\Term{twisted-constant-sheaf}{扭常值层} $T_\ell(A^0)^t$ 与 $T_\ell(A'^0)^t$ 的 $\overline K$ 值点群；这两个\Term{sheaf}{层}的\Term{special-fiber}{特殊纤维}分别是 $T_\ell(T_0)$ 与 $T_\ell(T'_0)$。于是得到\Term{canonical}{典范}\Term{isomorphism}{同构}
\begin{equation}\tag{9.2.7}
 \begin{cases}
 W\simeq T_\ell(T_0)(\overline k)\simeq M^\vee\otimes\mathbb Z_\ell(1),\\[4pt]
 W'\simeq T_\ell(T_0(\overline k))\simeq M'\otimes\mathbb Z_\ell(1),
 \end{cases}
\end{equation}
其中 $M$ 与 $M'$ 分别表示 $T_0$ 与 $T'_0$ 相对于 $\overline k$ 的\Term{character-group}{特征标群}：
\begin{equation}\tag{9.2.8}
 T_0=D(M),\qquad T'_0=D(M').
\end{equation}
所以，表达\Term{inertia-group}{惯性群} $I$ 在 $A_K$ 的 \Term{tate-module}{Tate 模} $U$ 上\Term{action}{作用}的 (9.2.6) 中的元素 $u$，\Term{solvable}{可解}释为 $(M'^\vee\otimes_{\mathbb Z}M^\vee)\otimes_{\mathbb Z}\mathbb Z_\ell$ 的一个\Term{canonical}{典范}元素。
\SourcePage{430}
这个元素还在 $\pi_0=\pi/I=\operatorname{Gal}(\overline k/k)$ 下\Term{invariant}{不变}，因为按定义它在\Term{transport-of-structure}{结构搬运}下\Term{invariant}{不变}：
\begin{equation}\tag{9.2.9}
 u\in(M'^\vee\otimes_{\mathbb Z}M^\vee\otimes_{\mathbb Z}\mathbb Z_\ell)^I
 \simeq\operatorname{Hom}_{\mathbb Z}(\underline M\otimes\underline M',\mathbb Z_\ell)
 \simeq\operatorname{Hom}_{\mathbb Z_\ell}(\underline M_\ell\otimes\underline M'_\ell,\mathbb Z_\ell).
\end{equation}
这正是我们要定义的 (9.1.2) 中的元素（仅将记号 $u_\ell$ 写成了 $u$）。

为在\Term{generic}{一般}情形定义 (9.1.2)，我们需要一些准备，下面两小节将展开这些准备。

\textbf{9.3。\Term{blended-extension}{混合扩张}。}设 $\mathcal C$ 是一个\Term{abelian-category}{阿贝尔范畴}，
\begin{equation}\tag{9.3.1}
 P,\ R,\ Q
\end{equation}
是 $\mathcal C$ 的三个\Term{object}{对象}。我们关心 $\mathcal C$ 中带有三步\Term{filtration}{滤过}
\begin{equation}\tag{*}
 X^0=X\supset X^1\supset X^2\supset X^3=0
\end{equation}
以及给定\Term{isomorphism}{同构}
\[
 X^0/X^1\xrightarrow{\sim}P,\qquad
 X^1/X^2\xrightarrow{\sim}R,\qquad
 X^2/X^3=X^2\xrightarrow{\sim}Q
\]
的\Term{object}{对象} $X$ 的分类。置
\[
 E=X^0/X^2,\qquad F=X^1/X^3=X^1,
\]
可见 $X$ 给出两个普通\Term{extension}{扩张}
\begin{equation}\tag{9.3.2}
 0\longrightarrow P\longrightarrow F\longrightarrow R\longrightarrow0,
 \qquad 0\longrightarrow R\longrightarrow E\longrightarrow Q\longrightarrow0,
\end{equation}
分别是 $R$ 被 $P$ 与 $Q$ 被 $R$ 的\Term{extension}{扩张}。
\SourcePage{431}
这里我们采取的观点是：除\Term{object}{对象} (9.3.1) 外，也从给定的\Term{extension}{扩张}结构 (9.3.2) 出发，考察相应的三步\Term{filtration}{滤过}结构 $X$。确切地说，我们把 $\mathcal C$ 的一个\Term{object}{对象} $X$ 连同三步\Term{filtration}{滤过} (*) 和给定\Term{isomorphism}{同构}
\begin{equation}\tag{9.3.3}
 E\xrightarrow{\sim}X^0/X^2,\qquad F\xrightarrow{\sim}X^1/X^3=X^1
\end{equation}
称为\Term{extension}{扩张} $E$ 被\Term{extension}{扩张} $F$ 的\Term{blended-extension}{混合扩张}（extension panachée），条件是这些\Term{isomorphism}{同构}与\Term{filtration}{滤过}相容，即分别把 $X^1/X^2$ 变为 $\operatorname{Im}(R\to E)$，把 $X^2$ 变为 $\operatorname{Im}(P,F)$；还要求由 (9.3.3) \Term{derive}{导出}的\Term{isomorphism}{同构} $R\simeq F/P$（$\simeq X^1/X^2$）等于由第一个\Term{exact-sequence}{正合列} (9.3.2) \Term{derive}{导出}的\Term{isomorphism}{同构}。固定 $P,Q,R$ 及\Term{extension}{扩张} $E,F$ 后，$E$ 被 $F$ 的\Term{blended-extension}{混合扩张}以显然的方式构成一个\Term{category}{范畴}，它实际上是一个\Term{groupoid}{群胚}（其中每个\Term{arrow}{箭头}都\Term{invertible}{可逆}）。一个\Term{blended-extension}{混合扩张} $X$ 的\Term{automorphism-group}{自同构群}\Term{canonical}{典范}地\Term{isomorphism}{同构}于群 $\operatorname{Hom}(Q,P)$，其\Term{isomorphism}{同构}为
\begin{equation}\tag{9.3.4}
 \operatorname{Hom}(Q,P)\xrightarrow{\sim}\operatorname{Aut}_{\mathrm{extpan}}(X),
 \qquad u\longmapsto\operatorname{id}_X+\overline u,
\end{equation}
其中 $\overline u$ 是 $X$ 的底层\Term{object}{对象}的\Term{endomorphism}{自同态}，由复合
\[
 X\longrightarrow Q\xrightarrow{u}P\longrightarrow X
\]
给出；两端\Term{arrow}{箭头}分别为 $X$ 的给定\Term{blended-extension}{混合扩张}结构所给出的\Term{canonical}{典范}\Term{epimorphism}{满态射}与\Term{canonical}{典范}\Term{monomorphism}{单态射}。

还可以进一步说明 $E$ 被 $F$ 的\Term{blended-extension}{混合扩张}所成\Term{category}{范畴} $\operatorname{EXTPAN}(E,F)$ 的结构：借助一个自然\Term{functor}{函子}（类似于普通\Term{extension}{扩张}的 Brauer 复合），把它看作\Term{extension}{扩张} $Q$ 被 $P$ 所成的“群\Term{category}{范畴}”（第 VII 讲 2.5）$\operatorname{EXT}(Q,P)$ 下的一个“\Term{torsor}{挠子}”。
\SourcePage{432}
这个\Term{functor}{函子}为
\begin{equation}\tag{9.3.5}
 \omega:\operatorname{EXT}(Q,P)\times\operatorname{EXTPAN}(E,F)
 \longrightarrow\operatorname{EXTPAN}(E,F),\qquad \omega(G,X)=G\wedge X,
\end{equation}
其性质是：对 $\operatorname{EXTPAN}(E,F)$ 的每个\Term{object}{对象} $X$，部分\Term{functor}{函子} $G\mapsto G\wedge X$ 都是从 $\operatorname{EXT}(Q,P)$ 到 $\operatorname{EXTPAN}(E,F)$ 的\Term{equivalence}{等价}。我们只概要描述\Term{functor}{函子} $\omega$。把 $E$ 被 $F$ 的\Term{blended-extension}{混合扩张} $X$ 看作 $Q$ 被 $F$ 的一个普通\Term{extension}{扩张}（仍记为 $X$，视为 $\mathcal C$ 中的\Term{object}{对象}）。对 $Q$ 被 $P$ 的每个\Term{extension}{扩张} $G$，以 $\overline G$ 表示借助包含 $P\to F$ 所确定的 $Q$ 被 $F$ 的\Term{extension}{扩张}。于是置
\begin{equation}\tag{9.3.6}
 G\wedge X=\omega(G,X)\overset{\mathrm{def}}=\overline G\wedge X,
\end{equation}
其中\Term{signature}{符号} $\wedge$ 表示 $Q$ 被 $F$ 的\Term{extension}{扩张}\Term{category}{范畴}中的 Brauer 复合。式 (9.3.6) 右端 $X'$ 的三步\Term{filtration}{滤过}显然可得，与 $X'$ 关联的 (9.3.3) 型的第二个\Term{isomorphism}{同构}也显然可得。为给 $X'$ 定义 (9.3.3) 型的第一个\Term{isomorphism}{同构}，注意：由给定\Term{morphism}{态射} $F\to R$ (9.3.2) \Term{derive}{导出}的\Term{functor}{函子}
\[
 Y\longmapsto\widetilde Y,\qquad
 \operatorname{EXT}(Q,F)\longrightarrow\operatorname{EXT}(Q,R)
\]
与 Brauer 复合相容，故将 $X'$ 变为一个\Term{canonical}{典范}地\Term{isomorphism}{同构}于 $\widetilde{\overline G}\wedge\widetilde X$ 的 $\widetilde X'$。第一个\Term{isomorphism}{同构} (9.3.3) 给出\Term{isomorphism}{同构} $\widetilde X\simeq F$；另一方面，$\widetilde{\overline G}$（由 $G$ 经复合 $P\to F\to R$ \Term{derive}{导出}的\Term{extension}{扩张}）显然\Term{canonical}{典范}地\Term{isomorphism}{同构}于\Term{trivial}{平凡}\Term{extension}{扩张}，因为所考察的复合 $P\to R$ 为零。由此得到\Term{canonical}{典范}\Term{isomorphism}{同构} $\widetilde X'\simeq F$，从而使 (9.3.6) 右端成为 $E$ 被 $F$ 的\Term{blended-extension}{混合扩张}。
\SourcePage{433}
我们把在 $\operatorname{EXT}(Q,P)\times\operatorname{EXTPAN}(E,F)$ 的\Term{arrow}{箭头}上定义 $\omega$、核实\Term{functor}{函子} (9.3.5) 确有上述性质（即对每个固定的\Term{blended-extension}{混合扩张} $X$，\Term{functor}{函子} $G\mapsto G\wedge X$ 是 $\operatorname{EXT}(Q,P)$ 与 $\operatorname{EXTPAN}(E,F)$ 之间的\Term{equivalence}{等价}），以及证明结合公式
\begin{equation}\tag{9.3.7}
 G_1\wedge(G_2\wedge X)\xrightarrow{\sim}(G_1\wedge G_2)\wedge X
\end{equation}
的工作留给读者；这里 $G_1\wedge G_2$ 表示 $\operatorname{EXT}(Q,P)$ 中的 Brauer 复合。

由上面的讨论立即得到下述命题的 a)、b) 两部分。
\begin{SGAresult}{命题}{9.3.8}
给定\Term{abelian-category}{阿贝尔范畴} $\mathcal C$ 中的两个\Term{extension}{扩张} $E,F$ (9.3.2)，考察 $E$ 被 $F$ 的\Term{blended-extension}{混合扩张}所成\Term{category}{范畴} $\operatorname{EXTPAN}(E,F)$。则有：

 a) $\operatorname{EXTPAN}(E,F)$ 是一个\Term{groupoid}{群胚}，其中每个\Term{object}{对象}的\Term{automorphism-group}{自同构群}均\Term{canonical}{典范}地\Term{isomorphism}{同构}于 $\operatorname{Ext}^0(Q,P)=\operatorname{Hom}(Q,P)$。

 b) $\operatorname{EXTPAN}(E,F)$ 中\Term{object}{对象}的\Term{isomorphism}{同构}类所成\Term{set}{集合} $\operatorname{Extpan}(E,F)$ 或者为空，或者自然地是群 $\operatorname{Ext}^1(Q,P)$ 下的一个\Term{torsor}{挠子}。

 c) 可以\Term{canonical}{典范}地定义一个类
\begin{equation}\tag{9.3.9}
 c(E,F)\in\operatorname{Ext}^2(Q,P),
\end{equation}
其为零是 $\operatorname{Extpan}(E,F)\ne\varnothing$，即存在 $E$ 被 $F$ 的\Term{blended-extension}{混合扩张}的充分必要条件。
\end{SGAresult}

\SourcePage{434}
\begin{proof}
还须证明 d)。注意，$E$ 被 $F$ 的一个\Term{blended-extension}{混合扩张}，也可以描述为 $E$ 被 $P$ 的一个\Term{extension}{扩张} $X$，附有 $R$ 被 $P$ 的\Term{extension}{扩张}之间的一个\Term{isomorphism}{同构}
\[
 X^1\xrightarrow{\sim}F,
\]
其中 $X^1$ 表示\Term{extension}{扩张} $X$ 经 (9.3.2) 给出的\Term{morphism}{态射} $R\to E$ 的像。因此，这样的\Term{blended-extension}{混合扩张}存在，当且仅当 $\operatorname{Ext}^1(E,P)$ 中有一个元素，其在 $\operatorname{Ext}^1(R,P)$ 中的像为\Term{extension}{扩张} $F$ 的类 $\xi$。考察与 (9.3.2) 的第二个\Term{exact-sequence}{正合列}关联的 $\operatorname{Ext}^i(-,P)$ 的\Term{exact-sequence}{正合列}
\[
 \cdots\longrightarrow\operatorname{Ext}^1(E,P)\longrightarrow
 \operatorname{Ext}^1(R,P)\xrightarrow{\partial}\operatorname{Ext}^2(R,Q)\longrightarrow\cdots,
\]
只须用下式定义 $c(E,F)$：
\begin{equation}\tag{9.3.10}
 c(E,F)=\partial\xi,\qquad \xi=\operatorname{cl}(F)\in\operatorname{Ext}^1(R,P).
\end{equation}
\end{proof}

现在设有一个\Term{exact-functor}{正合函子}
\begin{equation}\tag{9.3.11}
 r:\mathcal C\longrightarrow\mathcal C',
\end{equation}
其中 $\mathcal C'$ 是另一个\Term{abelian-category}{阿贝尔范畴}。以 $P',Q',R',E',F'$ 表示对 $P,Q,R,E,F$ 应用\Term{functor}{函子} $r$ 得到的 $\mathcal C'$ 中的\Term{object}{对象}；于是 $E'$ 是 $Q'$ 被 $R'$ 的\Term{extension}{扩张}，$F'$ 是 $R'$ 被 $P'$ 的\Term{extension}{扩张}。显然，下述\Term{functor}{函子}图在\Term{canonical}{典范}\Term{isomorphism}{同构}意义下交换：
\begin{equation}\tag{9.3.12}
\begin{tikzcd}[column sep=small]
 \operatorname{EXT}(Q,P)\times\operatorname{EXTPAN}(E,F) \arrow[r,"\omega"] \arrow[d]
 & \operatorname{EXTPAN}(E,F) \arrow[d]\\
 \operatorname{EXT}(Q',P')\times\operatorname{EXTPAN}(E',F') \arrow[r,"\omega'"]
 & \operatorname{EXTPAN}(E',F')
\end{tikzcd}
\end{equation}
\SourcePage{435}
这里竖直\Term{arrow}{箭头}是\Term{exact-functor}{正合函子} $r$ \Term{derive}{导出}的显然的\Term{functor}{函子}，$\omega'$ 则是将 $\mathcal C$ 换成 $\mathcal C'$ 后与 $\omega$ (9.3.5) 对应的\Term{arrow}{箭头}。由此得到 9.3.8 中所考察的结构与\Term{functor}{函子} $r$ 的类似相容性，具体陈述留给读者。

我们将使用这一\Term{commutativity}{交换性}的如下推论。设 $\operatorname{EXTPAN}(E,F)$ \Term{nonempty}{非空}，从中选择\Term{object}{对象} $X$。借助 $G'\mapsto G'\wedge X'$（其中 $X'=r(X)$），\Term{category}{范畴} $\operatorname{EXTPAN}(E',F')$ 与 $\operatorname{EXT}(Q',P')$ \Term{equivalence}{等价}。因此，只要知道对应的 $\operatorname{EXT}(Q',P')$ 中元素的\Term{isomorphism}{同构}类，就在\Term{isomorphism}{同构}意义下知道了 $\operatorname{EXTPAN}(E',F')$ 中的元素 $Y'$；也就是说，只须知道 $\operatorname{Ext}^1(Q',P')$ 中相应的元素，记为
\[
 c_X(Y')\in\operatorname{Ext}^1(Q',P').
\]
若将 $X$ 换成另一个\Term{object}{对象} $X_1=H\wedge X$，其中 $H$ 是 $Q$ 被 $P$ 的一个\Term{extension}{扩张}，则由结合公式 (9.3.7) 立即有
\[
 c_{X_1}(Y')=c_X(Y')-\operatorname{cl}(H').
\]
所以，固定 $Y'$ 而让 $X$ 变化时，所有 $c_X(Y')$ 恰好遍历一个类
\begin{equation}\tag{9.3.13}
 c(Y')\in\operatorname{Ext}^1(Q',P')/\operatorname{Im}\operatorname{Ext}^1(Q,P),
\end{equation}
这个类因而\Term{canonical}{典范}地与 $E'$ 被 $F'$ 的\Term{blended-extension}{混合扩张} $Y'$ 关联。按构造，其为零当且仅当\Term{blended-extension}{混合扩张} $Y'$ \Term{isomorphism}{同构}于 $E$ 被 $F$ 的某个\Term{blended-extension}{混合扩张} $Y$ 经 $r$ 变换所得的\Term{object}{对象}。

\SourcePage{436}
\textbf{9.4。}设 $S$ 是一个\Term{henselian}{Hensel}\Term{trait}{离散赋值环谱}，$\eta$ 为其\Term{generic-point}{一般点}；设 $n>0$ 是整数，$\Lambda=\mathbb Z/n\mathbb Z$，$\mathcal C$ 是 $S$ 上 $\Lambda$-\Term{module}{模}的\Term{category}{范畴}（取 $S$ 上的 fppf \Term{site}{景}，参见第 VIII 讲 0.2），$\mathcal C'$ 是 $\eta$ 上 $\Lambda$-\Term{module}{模}的\Term{category}{范畴}。设 $P,R,Q$ 是 $\mathcal C$ 中的\Term{object}{对象}，$E,F$ 分别是 $Q$ 被 $R$ 与 $R$ 被 $P$ 的\Term{extension}{扩张} (9.3.2)。假设 $Q$ 是 $\Lambda_S$ 上的\Term{finite-type}{有限型}\Term{locally-free}{局部自由}\Term{module}{模}，而 $P$ \Term{representable}{可表示}为 $S$ 上的一个\Term{finite}{有限}\Term{multiplicative-type}{乘法型}群（SGA 3 IX），从而具有\Term{form}{形式}
\[
 P=D(Q^*),
\]
其中 $D=\underline{\operatorname{Hom}}(-,\mathbb G_{mS})$ 表示\Term{cartier-duality}{Cartier 对偶}。考察 $P,R,Q,E,F$ 在 $\eta$ 上的\Term{restriction}{限制}，记为 $P_\eta,R_\eta,\ldots$；并设给定 $E_\eta$ 被 $F_\eta$ 的一个\Term{blended-extension}{混合扩张} (9.3) $X_\eta$。在这些条件下，我们将\Term{canonical}{典范}地定义一个\Term{pairing}{配对}
\begin{equation}\tag{9.4.1}
 c(X_\eta):Q\otimes_\Lambda Q^*\longrightarrow\Lambda_S
 \qquad(\Lambda=\mathbb Z/n\mathbb Z).
\end{equation}
为此，先设 $S$ \Term{strictly-local}{严格局部}。于是 $Q$ 与 $Q^*$ 都是\Term{constant-group}{常值}的，记为 $Q=Q_{0S}$，$Q^*=Q^*_{0S}$；给出 (9.4.1)，\Term{equivalence}{等价}于给出一个元素
\begin{equation}\tag{9.4.2}
 c(X_\eta)\in Q_0^\vee\otimes Q_0^{*\vee},
\end{equation}
这里 $\vee$ 表示通常的\Term{dual}{对偶} $\operatorname{Hom}(-,\Lambda)$。我们应用 9.3 的讨论，取\Term{restriction}{限制}\Term{functor}{函子}为 $r:\mathcal C\to\mathcal C'$。先注意，存在 $E$ 被 $F$ 的一个\Term{blended-extension}{混合扩张}：事实上，这种\Term{extension}{扩张}存在的\Term{obstruction}{障碍}属于 $\operatorname{Ext}^2(Q,P)$（9.3.8 c)），只须证明
\[
 \operatorname{Ext}^2(Q,P)=0.
\]
\SourcePage{437}
这是下述\Term{lemma}{引理}的一个特例。
\begin{SGAresult}{\Term{lemma}{引理}}{9.4.3}
在上述条件下，若 $S$ \Term{strictly-local}{严格局部}，则有\Term{isomorphism}{同构}
\begin{align*}
 \text{a)}\quad
 \operatorname{Ext}^i_\Lambda(Q,P)&\simeq H^i(S,Q^\vee\otimes P),\\
 \operatorname{Ext}^i_\Lambda(Q_\eta,P_\eta)&\simeq H^i(\eta,Q^\vee\otimes P).
\end{align*}
另一方面，有\Term{canonical}{典范}\Term{isomorphism}{同构}
\[
 \text{b)}\qquad Q^\vee\otimes P\simeq D(Q\otimes Q^*).
\]
最后，对 $S$ 上的每个\Term{finite}{有限}\Term{constant-group}{常值群} $H=H_{0S}$，有
\[
 \text{c)}\qquad H^i(S,D(H))=H^i(\eta,D(H)_\eta)=0\qquad(i\geq2),
\]
以及\Term{canonical}{典范}\Term{isomorphism}{同构}
\begin{align*}
 \text{d)}\quad H^1(S,D(H))&\simeq\mathbb G_m(S)\otimes H_0^\vee
   =V^*\otimes H_0^\vee,\\
 H^1(\eta,D(H)_\eta)&\simeq\mathbb G_m(\eta)\otimes H_0^\vee
   =K^*\otimes H_0^\vee,
\end{align*}
这里\Term{signature}{符号} $\vee$ 表示 $\operatorname{Hom}(-,\mathbb Q/\mathbb Z)$。
\end{SGAresult}
\begin{proof}
公式 a)、b) 直接来自 $Q$ 是 $\Lambda_S$ 上的\Term{finite-type}{有限型}\Term{locally-free}{局部自由}\Term{module}{模}这一事实。为证明 c)，显然可以化归到 $H_0=\mathbb Z/m\mathbb Z$，即 $D(H)=\mu_m$ 的情形；于是，与 Kummer \Term{exact-sequence}{正合列}
\[
 0\longrightarrow\mu_m\longrightarrow\mathbb G_m
 \xrightarrow{m.\operatorname{id}}\mathbb G_m\longrightarrow0
\]
关联的\Term{cohomology}{上同调}\Term{exact-sequence}{正合列}，把我们归结为证明公式
\[
 H^i(S,\mathbb G_m)=H^i(\eta,\mathbb G_m)=0\qquad(i\geq1).
\]
这里的 $H^i$ 等于在\Term{etale}{平展}\Term{site}{景}上计算的 $H^i_{\text{ét}}$ [6, th.~11.7]；对 $S$，因其\Term{strictly-local}{严格局部}而为零。它们对 $\eta$ 也为零，则是一个熟知的事实 [6, th.~1.1]。

\SourcePage{438}
还须描述 d) 中的\Term{isomorphism}{同构}。事实上，在任意底 $U$ 上，都有一个\Term{canonical}{典范}\Term{homomorphism}{同态}
\begin{equation}\tag{9.4.3.1}
 \mathbb G_m(U)\otimes H_0^\vee\longrightarrow H^1(U,D(H_{0U})).
\end{equation}
若 $H_0$ 被整数 $m\geq0$ \Term{annihilate}{零化}，就可以借助与 Kummer \Term{exact-sequence}{正合列}关联的上边缘\Term{homomorphism}{同态} $\mathbb G_m(U)\xrightarrow{\partial}H^1(U,\mu_m)$，将它描述为下述复合：
\[
 \begin{aligned}
 \mathbb G_m(U)\otimes H_0^\vee
 &\longrightarrow H^1(U,\mu_m)\otimes H_0^\vee
 \longrightarrow H^1(U,\mu_m\otimes H_0^\vee)\\
 &\longrightarrow H^1(U,D(H_{0S})),
 \end{aligned}
\]
其中第二个\Term{arrow}{箭头}来自显然的\Term{homomorphism}{同态}
\[
 H_0^\vee\longrightarrow\operatorname{Hom}(\mu_m,\mu_m\otimes H_0^\vee),
\]
第三个\Term{arrow}{箭头}则是 b) 的一个特例（把 $n$ 换成 $m$）。我们把核实 (9.4.3.1) 不依赖于满足 $mH_0=0$ 的整数 $m$ 的选择这件事留给读者。

另一方面，我断言，\Term{homomorphism}{同态} (9.4.3.1) 总是\Term{injective}{单射}的，而当 $\operatorname{Pic}(U)=0$ 时，它的\Term{cokernel}{余核}为零。（实际上，这个\Term{cokernel}{余核}总\Term{isomorphism}{同构}于 $\operatorname{Hom}(H_0,\operatorname{Pic}(U))$，但这里不必用到。）为证明这一点，仍可化归到 $H_0$ 形如 $\mathbb Z/n\mathbb Z$ 的情形，此时断言立即来自上面的 Kummer \Term{exact-sequence}{正合列} (*) 所关联的\Term{cohomology}{上同调}\Term{exact-sequence}{正合列}。这就证明了公式 d)。
\end{proof}
立即得到：
\begin{SGAresult}{推论}{9.4.4}
沿用 9.4.3 的记号，有\Term{canonical}{典范}\Term{isomorphism}{同构}
\[
 \text{e)}\qquad
 H^1(\eta,D(H)_\eta)/\operatorname{Im}H^1(S,D(H))\simeq H_0^\vee,
\]
进而（结合 a)、b)）有\Term{canonical}{典范}\Term{isomorphism}{同构}
\begin{equation}\tag{9.4.5}
 \operatorname{Ext}^1_\Lambda(Q_\eta,P_\eta)/\operatorname{Im}\operatorname{Ext}^1_\Lambda(Q,P)
 \simeq Q_0^\vee\otimes Q_0^{*\vee}.
\end{equation}
\end{SGAresult}

\SourcePage{439}
回到 $S$ \Term{strictly-local}{严格局部}时 (9.4.2) 的定义。关系 $\operatorname{Ext}^2(Q,P)=0$ 使我们能够定义 (9.3.13) 中的\Term{canonical}{典范}元素，其中取给定的 $E_\eta$ 被 $F_\eta$ 的\Term{blended-extension}{混合扩张} $X_\eta$ 为 $Y'$。公式 (9.4.5) 将这个元素 $c(X_\eta)$ 视为 $Q_0^\vee\otimes Q_0^{*\vee}$ 中的元素，从而明确了 (9.4.2) 的定义。由 9.3 的\Term{generic}{一般}讨论还可知，这个元素可以解释为把 $X_\eta$ 延拓成 $E$ 被 $F$ 的一个\Term{blended-extension}{混合扩张} $X$ 的\Term{obstruction}{障碍}。

\textbf{9.4.6。}当 $S$ 为\Term{henselian}{Hensel}时，只须从 $S$ 的\Term{strict-henselization}{严格 Hensel 化}作\Term{galois-descent}{Galois 下降}，便立即把 (9.4.1) 的定义归结为\Term{strictly-local}{严格局部}情形。

\begin{SGAresult}{命题}{9.4.7}
数据 $P,Q,R,E,F$ 与上面相同，但不必假设\Term{trait}{离散赋值环谱} $S$ 为\Term{henselian}{Hensel}。考察\Term{functor}{函子}
\begin{equation}\tag{9.4.7.1}
 X\longmapsto X_\eta,
\end{equation}
从 $E$ 被 $F$ 的\Term{blended-extension}{混合扩张}\Term{category}{范畴}到 $E_\eta$ 被 $F_\eta$ 的\Term{blended-extension}{混合扩张}\Term{category}{范畴}。对 $E_\eta$ 被 $F_\eta$ 的每个\Term{blended-extension}{混合扩张} $X_\eta$，考察\Term{pairing}{配对}
\begin{equation}\tag{9.4.7.2}
 c_0(X_\eta):Q_0\otimes Q_0^*\longrightarrow\Lambda_S,
\end{equation}
它由在 $S$ 的\Term{henselization}{Hensel 化} $S^h$ 上定义的 (9.4.1) 型\Term{pairing}{配对}诱导。则\Term{functor}{函子} (9.4.7.1) \Term{fully-faithful}{全忠实}；$E_\eta$ 被 $F_\eta$ 的一个\Term{blended-extension}{混合扩张} $X_\eta$ 属于这个\Term{functor}{函子}的\Term{essential-image}{本质像}，当且仅当\Term{pairing}{配对} (9.4.7.2) 为零。
\end{SGAresult}

\SourcePage{440}
\textbf{9.4.7.3。}若只假设 $S$ 是一个一维\Term{noetherian}{诺特}\Term{regular}{正则}\Term{connected}{连通}\Term{scheme}{概形}，这个陈述仍有意义且仍然成立：此时不能只考虑一个\Term{pairing}{配对} (9.4.7.2)，而须对 $S$ 的每个\Term{closed-point}{闭点} $s$ 考虑一个\Term{pairing}{配对} $Q_s\otimes Q_s^*\to\Lambda_s$。以这个较\Term{generic}{一般}的\Term{form}{形式}来证明 9.4.7 最为方便。
\begin{proof}
 a) 对每个\Term{blended-extension}{混合扩张} $X$，映射
\[
 \operatorname{End}(X)\longrightarrow\operatorname{End}(X_\eta)
\]
是\Term{bijection}{双射}。由于两端分别是 $\underline{\operatorname{Hom}}(P,Q)=D(Q\otimes Q^*)$ 在 $S$ 与 $\eta$ 上的\Term{section}{截面}，这个断言来自 $S$ \Term{normal}{正规}而 $D(Q\otimes Q^*)$ 在 $S$ 上\Term{finite}{有限}这一事实。由此立即可得：若 $X,Y$ 是两个\Term{blended-extension}{混合扩张}，并且已经知道它们\Term{isomorphism}{同构}，则
\[
 \operatorname{Hom}(X,Y)\longrightarrow\operatorname{Hom}(X_\eta,Y_\eta)
\]
是\Term{bijection}{双射}。

 b) 若 $X,Y$ 是两个\Term{blended-extension}{混合扩张}，则每个\Term{isomorphism}{同构} $X_\eta\xrightarrow{\sim}Y_\eta$ 都来自一个\Term{isomorphism}{同构} $X\to Y$。由 a)，这个问题对于\Term{etale}{平展}\Term{topology}{拓扑}是\Term{local-scheme}{局部}的，因此可以化归到 $S$ \Term{strictly-local}{严格局部}的情形。但在这种情形，9.4.3 表明映射
\[
 \operatorname{Ext}^1_\Lambda(P,Q)\longrightarrow\operatorname{Ext}^1_\Lambda(P_\eta,Q_\eta)
\]
可视为映射
\[
 V^*\otimes H_0^\vee\longrightarrow K^*\otimes H_0^\vee,
\]
\SourcePage{441}
后者是\Term{injective}{单射}（因为\Term{exact-sequence}{正合列} $0\to V^*\to K^*\to\mathbb Z\to0$ \Term{split}{分裂}），故前者也是\Term{injective}{单射}；由此（借助 9.3 中的\Term{torsor}{挠子}）可知，若 $X_\eta$ 与 $Y_\eta$ \Term{isomorphism}{同构}，则 $X$ 与 $Y$ 也\Term{isomorphism}{同构}。再由 a) 即得所需结论。

显然，a)、b) 蕴含\Term{functor}{函子} $X\mapsto X_\eta$ \Term{fully-faithful}{全忠实}。

 c) 设 $X_\eta$ 是 $\eta$ 上的一个\Term{blended-extension}{混合扩张}，我们证明：它可延拓到 $S$ 当且仅当各\Term{pairing}{配对} (9.4.6.2) 为零。由已经证明的唯一性，问题仍然对于\Term{etale}{平展}\Term{topology}{拓扑}是\Term{local-scheme}{局部}的，于是又归结到 $S$ \Term{strictly-local}{严格局部}的情形。此时结论按构造成立，正如 9.4.6 前已经指出的。

至此，9.4.7 证毕。
\end{proof}

\textbf{9.5。}现在固定一个\Term{prime-number}{素数} $\ell$、$S$ 上的 Barsotti--Tate \Term{pro-ell-group}{pro-$\ell$ 群}
\begin{equation}\tag{9.5.1}
 P=(P(n))_{n\geq0},\qquad R=(R(n))_{n\geq0},\qquad Q=(Q(n))_{n\geq0},
\end{equation}
以及 Barsotti--Tate \Term{pro-ell-group}{pro-$\ell$ 群}的\Term{extension}{扩张}
\begin{equation}\tag{9.5.2}
 0\longrightarrow P\longrightarrow F\longrightarrow R\longrightarrow0,
 \qquad0\longrightarrow R\longrightarrow E\longrightarrow Q\longrightarrow0.
\end{equation}
于是，对每个整数 $n\geq0$，它们定义 $\Lambda_n$-\Term{module}{模}的\Term{extension}{扩张}，其中 $\Lambda_n=\mathbb Z/\ell^{n+1}\mathbb Z$：
\begin{equation}\tag{9.5.3}
 \begin{gathered}
 0\longrightarrow P(n)\longrightarrow F(n)\longrightarrow R(n)\longrightarrow0,\\
 0\longrightarrow R(n)\longrightarrow E(n)\longrightarrow Q(n)\longrightarrow0.
 \end{gathered}
\end{equation}
我们假设 $Q$ 为\Term{etale}{平展}型，而 $P$ 为\Term{multiplicative-type}{乘法型}。
\SourcePage{442}
最后，设 $\eta$ 上给定一个 Barsotti--Tate \Term{pro-ell-group}{pro-$\ell$ 群} $X_\eta$，它是 $E_\eta$ 被 $F_\eta$ 的\Term{blended-extension}{混合扩张}；其含义是：它（作为 Barsotti--Tate \Term{pro-ell-group}{pro-$\ell$ 群}）带有三步\Term{filtration}{滤过}，由此产生的两个部分\Term{extension}{扩张}分别附有与 $E_\eta$、$F_\eta$ 的\Term{isomorphism}{同构}。因此，对每个 $n$，$X_\eta(n)$ 带有 $E(n)_\eta$ 被 $F(n)_\eta$ 的\Term{blended-extension}{混合扩张}结构。于是可将它对应到一个 (9.4.1) 型\Term{pairing}{配对}：
\begin{equation}\tag{9.5.4}
 c(X(n)_\eta):Q(n)\otimes Q^*(n)\longrightarrow\Lambda_n
 =\mathbb Z/\ell^{n+1}\mathbb Z,
\end{equation}
其中 $Q^*(n)$ 是 $S$ 上由下式定义的\Term{finite}{有限}\Term{etale}{平展}群：
\[
 P(n)\simeq D(Q^*(n)),\qquad\text{即}\quad Q^*(n)=D(P(n)),
\]
$D$ 仍表示\Term{cartier-duality}{Cartier 对偶}。各 $Q^*(n)$ 以显然的方式构成一个\Term{pro-ell-group}{pro-$\ell$ 群}，记为 $Q^*$。当 $n$ 变化时，各\Term{homomorphism}{同态} (9.5.4) 定义一个\Term{projective-system}{投射系}的\Term{homomorphism}{同态}，即 $S$ 上 $\ell$ \Term{adic}{进}\Term{twisted-constant-sheaf}{扭常值层}的\Term{homomorphism}{同态}
\begin{equation}\tag{9.5.5}
 c(X_\eta):Q\otimes Q^*\longrightarrow\mathbb Z_{\ell S};
\end{equation}
我们把这件事的核实留给读者。

顺便指出以下结果（后文不会用到它）：
\begin{SGAresult}{命题}{9.5.6}
数据 $P,Q,R,E,F$ 与上面相同，但不必假设\Term{trait}{离散赋值环谱} $S$ 为\Term{henselian}{Hensel}。考察\Term{functor}{函子}
\begin{equation}\tag{9.5.6.1}
 X\longmapsto X_\eta,
\end{equation}
从 $E$ 被 $F$ 的 Barsotti--Tate \Term{blended-extension}{混合扩张}所成\Term{category}{范畴}，到 $E_\eta$ 被 $F_\eta$ 的\Term{blended-extension}{混合扩张}所成\Term{category}{范畴}。
\SourcePage{443}
对每个作为 $E_\eta$ 被 $F_\eta$ 的\Term{blended-extension}{混合扩张}的 Barsotti--Tate 群 $X_\eta$，考察相应\Term{pairing}{配对}
\begin{equation}\tag{9.5.6.2}
 Q_0\otimes Q_0^*\longrightarrow\mathbb Z_{\ell S},
\end{equation}
它由 $S$ 的\Term{henselization}{Hensel 化}上定义的 (9.5.5) 型\Term{pairing}{配对}诱导。则\Term{functor}{函子} (9.5.6.1) \Term{fully-faithful}{全忠实}，其\Term{essential-image}{本质像}恰由那些相应\Term{pairing}{配对} (9.5.6.2) 为零的、作为 $E_\eta$ 被 $F_\eta$ 的\Term{blended-extension}{混合扩张}的 Barsotti--Tate 群 $X_\eta$ 构成。此外，当 $\operatorname{char}K=0$（从而可应用 Tate \Term{theorem}{定理} [33, th.~4]）时，$X_\eta$ 属于\Term{essential-image}{本质像}，当且仅当 $X_\eta$ 可延拓为 $S$ 上的一个 Barsotti--Tate 群。
\end{SGAresult}

\textbf{9.5.6.3。}若只假设 $S$ 是一个一维\Term{noetherian}{诺特}\Term{regular}{正则}\Term{connected}{连通}\Term{scheme}{概形}，上述陈述也可推广；我们把这件事留给读者。
\begin{proof}
把 9.4.6 应用于所考察的 Barsotti--Tate 群的各个\Term{set-component}{分量}，立即得到 9.5.6 的第一个断言。还须证明：若作为 $E_\eta$ 被 $F_\eta$ 的\Term{blended-extension}{混合扩张}的 Barsotti--Tate 群 $X_\eta$ 可以延拓为 $S$ 上的一个 Barsotti--Tate 群 $X$，则它也可作为\Term{blended-extension}{混合扩张}延拓。根据 Tate \Term{theorem}{定理} [33]，\Term{homomorphism}{同态} $X_\eta\to Q_\eta$ 唯一地延拓成\Term{homomorphism}{同态}
\[
 X\longrightarrow Q.
\]
\SourcePage{444}
对每个整数 $n$，相应的\Term{homomorphism}{同态} $X(n)\to Q(n)$ 在\Term{generic-fiber}{泛纤维}上\Term{surjective}{满射}，因而\Term{surjective}{满射}；由于 $Q(n)$ 在 $\eta$ 上\Term{etale}{平展}，故 $X(n)\to Q(n)$ \Term{flat}{平坦}，所以它是一个\Term{epimorphism}{满态射}。因此，$X\to Q$ 是逐\Term{set-component}{分量}的\Term{epimorphism}{满态射}。于是，若以 $\overline F$ 表示 $X\to Q$ 的\Term{kernel}{核}，则 $\overline F$ 是一个 Barsotti--Tate 群，而 $X$ 是 $Q$ 被 $\overline F$ 的\Term{extension}{扩张}。另外，由 Tate \Term{theorem}{定理}，\Term{isomorphism}{同构} $F_\eta\xrightarrow{\sim}\overline F_\eta$ 延拓成一个\Term{isomorphism}{同构} $F\simeq\overline F$。所以 $X$ 成为 $Q$ 被 $F$ 的一个\Term{extension}{扩张}。再考察\Term{quotient}{商} $\overline E=X/P$，由 Tate \Term{theorem}{定理}又可得一个\Term{isomorphism}{同构} $E\simeq\overline E$。因此，$X$ 是延拓给定\Term{blended-extension}{混合扩张} $X_\eta$ 的 $E$ 被 $F$ 的\Term{blended-extension}{混合扩张}，证毕。
\end{proof}

\textbf{注记 9.5.6.4。}a) 当然，9.5.6 的最后一个断言理应对 $K$ 不加\Term{restriction}{限制}也成立，因为所引用的 Tate \Term{theorem}{定理}理应也是如此。

b) 可给出\Term{good-reduction}{好约化}判别法 5.10 的另一个证明：先如 5.10 中那样化归到 $A_K$ 具有\Term{semistable-reduction}{半稳定约化}的情形，然后将 9.5.6 用于作为\Term{blended-extension}{混合扩张}的 $T_\ell(A_K)$（参见 9.6），得到 $u_\ell=0$，再使用下文的 11.6.2 a)，便有 $T_\ell(A_K)^t=0$。

\textbf{9.5.8。}当 $\ell\ne p$ 时，$S$ 上（相应地，$K$ 上）所考察的 Barsotti--Tate 群，可以简单地解释为 $S$ 上的 $\ell$ \Term{adic}{进}\Term{twisted-constant-sheaf}{扭常值层}；\Term{equivalence}{等价}地，可以解释为 $\pi_0=\operatorname{Gal}(\overline k/k)$（相应地，$\pi=\operatorname{Gal}(\overline K/K)$）在\Term{finite-type}{有限型}\Term{free-module}{自由 $\mathbb Z_\ell$-模}上的\Term{representation}{表示}。
\SourcePage{445}
此时可以看出，\Term{homomorphism}{同态} (9.5.6) 正是用 9.2 的方法定义的、表达\Term{inertia-group}{惯性子群} $I\subset\pi$ 在 $X_\eta(\overline K)$ 上\Term{action}{作用}的\Term{canonical}{典范}\Term{homomorphism}{同态}。为核实这一点，立即可以化归到 9.4 的情境，并假设其中 $n$ 等于 $\ell$ 的一个幂，而 $R$ 是\Term{locally-constant}{局部常值} $\Lambda$-\Term{module}{模}（与 $Q,P$ 一样），从而 $E,F,X_\eta$ 也同样如此。两个定义相同的核实留给读者。

\textbf{9.6。}现在已能定义所预告的\Term{canonical}{典范}\Term{pairing}{配对} (9.1.2)，无需对 $\ell$ 加\Term{restriction}{限制}，并使得在 $\ell\ne p$ 时这个定义与 9.2 给出的定义一致。为此，将 9.5 的讨论用于如下情境：
\begin{equation}\tag{9.6.1}
 \begin{cases}
 P=T_\ell(A^0)^t=\underline M_\ell'{}^\vee(1),
 \quad Q=\underline M_\ell,
 \quad R=T_\ell(A^0)^f/T_\ell(A^0)^t,\\[4pt]
 F=T_\ell(A^0)^f,\qquad E=D(F')=D(T_\ell(A'^0)^f).
 \end{cases}
\end{equation}
在最后一个公式中，$D$ 表示 Barsotti--Tate \Term{pro-ell-group}{pro-$\ell$ 群}意义下的\Term{cartier-dual}{Cartier 对偶}，即 $D(Y)$ 是由各 $D(Y(n))$ 以及通常的过渡\Term{morphism}{态射}构成的\Term{projective-system}{投射系}。$F$ 作为 $R$ 被 $P$ 的\Term{extension}{扩张}的结构是显然的。类似地，$F'=T_\ell(A'^0)^f$ 具有 $R'$ 被 $P'=T_\ell(A'^0)^t=\underline M_\ell^\vee(1)$ 的\Term{extension}{扩张}结构。应用\Term{cartier-duality}{Cartier 对偶}，可得 $E=D(F')$ 带有 $D(P')=\underline M_\ell$ 被 $D(R')$ 的\Term{extension}{扩张}结构。但根据 5.5（在 7.4.6 中证明），$D(R')$ \Term{canonical}{典范}地\Term{isomorphism}{同构}于 $R$，这就明确了 $E$ 作为 $Q$ 被 $R$ 的\Term{extension}{扩张}的结构。

\SourcePage{446}
作为 $E_\eta$ 被 $F_\eta$ 的\Term{blended-extension}{混合扩张} $X_\eta$，取
\begin{equation}\tag{9.6.2}
 X_\eta=T_\ell(A_K)=X_\eta^0\supset X_\eta^1=T_\ell(A_K)^f
 \supset X_\eta^2=T_\ell(A_K)^t\supset0.
\end{equation}
$R_\eta$ 被 $P_\eta$ 的\Term{extension}{扩张}之间的\Term{isomorphism}{同构}
\[
 X_\eta^1\xrightarrow{\sim}F_\eta
\]
是显然的。由\Term{orthogonality-theorem}{正交性定理} 5.2，$Q$ 被 $R$ 的\Term{extension}{扩张}之间的\Term{isomorphism}{同构}
\[
 X_\eta/X_\eta^2\xrightarrow{\sim}F_\eta=D(F'_\eta)
\]
也是显然的。

这样便具备了定义一个\Term{pairing}{配对} (9.5.6) 所需的数据；这里由于 $Q^*=\underline M'_\ell$，它就具有 (9.1.2) 的\Term{form}{形式}。这个定义与 9.2 的定义相容，来自 9.5.8。另一方面，5.8 的讨论表明，若 Tate \Term{theorem}{定理} [33] 成立（故至少在 $\operatorname{char}K=0$ 时成立），所考察的\Term{pairing}{配对}可仅用 $S$ 上的 Barsotti--Tate \Term{pro-ell-group}{pro-$\ell$ 群} $X_\eta=T_\ell(A_K)$ 来描述（因为它的三步\Term{filtration}{滤过}，以及延拓 $X_\eta^1$ 与 $X_\eta/X_\eta^2$ 的\Term{extension}{扩张} $F,E$，均可仅由\Term{pro-ell-group}{pro-$\ell$ 群} $X=T_\ell(A_K)$ \Term{canonical}{典范}地定义）。

\textbf{注记 9.7。}现在不再要求 $A_K$ 在\Term{henselian}{Hensel}\Term{trait}{离散赋值环谱} $S$ 上一定具有\Term{semistable-reduction}{半稳定约化}。设 $K'$ 是 $K$ 的一个\Term{finite}{有限}\Term{galois-extension}{Galois 扩张}，使 $A_{K'}$ 在 $S$ 于 $K'$ 中的\Term{normalization}{正规化} $S'$ 上具有\Term{semistable-reduction}{半稳定约化}（3.6）。沿用 (8.1.8) 的记号，由 $A_{K'}$ 的\Term{monodromy-pairing}{单值性配对}作\Term{galois-descent}{Galois 下降}，可定义一个\Term{pairing}{配对}
\SourcePage{447}
\begin{equation}\tag{*}
 v_\ell^{K'}:\underline M_K\otimes\underline M'_K\longrightarrow\mathbb Z_{\ell K}.
\end{equation}
它先验地依赖于 $K'$ 的选择，但由下文 (10.4.3) 可知，由
\[
 u_\ell=\frac1{e(S'/S)}v_\ell^{K'}
\]
给出的\Term{pairing}{配对}
\begin{equation}\tag{9.7.1}
 u_\ell:\underline M_K\otimes\underline M'_K\longrightarrow\mathbb Q_{\ell K}
\end{equation}
不依赖于这一选择，其中 $e(S'/S)=\operatorname{long}(V'/\mathfrak m V')$。仍可把 (9.7.1) 称为 $K$ 上的\Term{abelian-scheme}{阿贝尔概形} $A_K$ 对所考察的\Term{prime-number}{素数} $\ell$ 的\Term{monodromy-pairing}{单值性配对}，因为当 $A_K$ 本身具有\Term{semistable-reduction}{半稳定约化}时，它显然与\Term{monodromy-pairing}{单值性配对} (9.1.2) 在 $\eta$ 上的\Term{restriction}{限制}一致。由下文 10.4 可知，这个\Term{pairing}{配对}实际上总取值于 $\mathbb Q_K$，并且与 $\ell$ 无关。

可以问，它是否取值于 $\mathbb Z_K$；\Term{equivalence}{等价}地，是否对每个\Term{prime-number}{素数} $\ell$，(9.7.1) 都取值于 $\mathbb Z_{\ell K}$。至少当 $A_K$ 是\Term{elliptic-curve}{椭圆曲线}时确实如此。此时唯一须核实的，是 $\underline M_K\simeq\underline M'_K$ 的\Term{rank}{秩}为 $1$ 的情形；于是 $\underline M_K\otimes\underline M'_K$ \Term{canonical}{典范}地\Term{isomorphism}{同构}于 $\mathbb Z_K$，而\Term{monodromy-pairing}{单值性配对} (9.7.1) 便可简单地解释为一个普通有理数。由 Néron 的结果 [19] 可得，这个数等于 $v(j)$，其中 $j$ 是 $A_K$ 的绝对\Term{invariants}{不变量}，$v$ 是 $K$ 的规范化\Term{valuation}{赋值}。
